% UBT-AI-PROVENANCE-BEGIN
% schema: ubt-ai-provenance/v1
% tier: C_working
% ai_assistance: disclosed
% human_review: risk-based
% editorial_responsibility: Ing. David Jaroš
% policy: ../../../AI_PROVENANCE.md
% notice: Working material; exhaustive human review is not claimed.
% UBT-AI-PROVENANCE-END

\chapter{Formální důkazy a lemmata}

Tento dodatek shrnuje tři krátké argumenty, které se v učebnici opakovaně
používají. Jsou uvedeny zde, aby hlavní kapitoly zůstaly čitelné a současně
nepřebíraly standardní algebru jako vnější předpoklad.

\section{Centrální antikomutátor na Lorentzově řezu}
Nechť
\[
 E_\mu=i e_\mu{}^0\mathbf1+e_\mu{}^k\mathbf e_k,
 \qquad
 E_\mu^\sharp=i e_\mu{}^0\mathbf1-e_\mu{}^k\mathbf e_k.
\]
Při použití vztahu
\[
 \mathbf e_j\mathbf e_k=-\delta_{jk}\mathbf1
 +\epsilon_{jk\ell}\mathbf e_\ell
\]
se antisymetrické členy odpovídající vektorovému součinu po symetrizaci
zruší. Proto
\[
 \frac12\left(E_\mu^\sharp E_\nu+E_\nu^\sharp E_\mu\right)
 =\left(-e_\mu{}^0e_\nu{}^0+\sum_{k=1}^3e_\mu{}^ke_\nu{}^k\right)\mathbf1.
\]
Koeficientem je tetrádová metrika
$g_{\mu\nu}=e_\mu{}^ae_\nu{}^b\eta_{ab}$. Není třeba stopa ani projekce
na reálnou část.

\section{Hodnost zobrazení tetrády na metriku}
Pro
\[
 g_{\mu\nu}=e_\mu{}^ae_\nu{}^b\eta_{ab}
\]
je diferenciál
\[
 \delta g_{\mu\nu}
 =\eta_{ab}\left(\delta e_\mu{}^a e_\nu{}^b
 +e_\mu{}^a\delta e_\nu{}^b\right).
\]
U nedegenerované tetrády vznikne každý symetrický tenzor $h_{\mu\nu}$ volbou
\[
 \delta e_\mu{}^a=\frac12h_{\mu\rho}e^{\rho a}.
\]
Obraz má tedy rozměr deset. Jádro tvoří šest infinitezimálních
Lorentzových transformací
$\delta e_\mu{}^a=\lambda^a{}_b e_\mu{}^b$ s
$\lambda_{ab}=-\lambda_{ba}$.

\section{Proč variace není derivace}
Časoprostorová derivace a variace v prostoru polí působí v odlišných
směrech:
\[
 D_\mu\Theta=\partial_\mu\Theta+\rho_*(\Omega_\mu)\Theta,
 \qquad
 \Theta_\varepsilon=\Theta+\varepsilon\,\delta\Theta.
\]
Závisí-li konexe na poli, dává řetězové pravidlo
\[
 \delta(D_\mu\Theta)
 =D_\mu(\delta\Theta)+\rho_*(\delta\Omega_\mu)\Theta.
\]
Vynechání druhého členu mění kompozitní Eulerův--Lagrangeův systém.
Proto nelze výpočet s pevnou konexí automaticky znovu použít po dosazení
kompozitní konexe.
